\documentclass[11pt]{article}
\usepackage{amsmath,amssymb}
\usepackage[margin=1in]{geometry}
\begin{document}

\section{Cooperative diffusivity from the displacement transients}

In one dimension the solvent mass balance, Darcy's law and the Hookean network stress $\sigma'_{zz} = M\,\partial u_z/\partial z$ combine to
\begin{equation}
\frac{\partial u_z}{\partial t} = q(t) + \mathcal{D}_c \frac{\partial^2 u_z}{\partial z^2}, \qquad \mathcal{D}_c \equiv \kappa M,
\label{eq:upde}
\end{equation}
where $q(t) = \phi_s v_s + \phi_p v_p$ is the total volumetric flux, uniform in $z$. The strain $\varepsilon \equiv -\partial u_z/\partial z$ then obeys
\begin{equation}
\frac{\partial \varepsilon}{\partial t} = \mathcal{D}_c \frac{\partial^2 \varepsilon}{\partial z^2}
\label{eq:epde}
\end{equation}
on $0 \le z \le L$, with $L$ the thickness of the configuration the modes are written on and $\zeta \equiv z/L$. Each protocol below fixes the boundary and initial conditions of (\ref{eq:epde}); the fitted quantity is the decay rate of the first mode, $\tau_1^{-1}$, so $\mathcal{D}_c \propto L^2$ for a given $\tau_1$.

\subsection{Compression, loading}

After the ramp both plates are held, $u_z(0,t) = 0$ and $u_z(L,t) = -\Delta L$, and both faces drain into the same bath, so the strain is periodic:
\begin{equation}
\varepsilon(0,t) = \varepsilon(L,t), \qquad \frac{\partial\varepsilon}{\partial z}\Big|_{0} = \frac{\partial\varepsilon}{\partial z}\Big|_{L}.
\label{eq:holdbc}
\end{equation}
The solution of (\ref{eq:epde}) and (\ref{eq:holdbc}) is
\begin{equation}
\varepsilon(\zeta,t) = \frac{\Delta L}{L} + \sum_{m=1}^{\infty} \left[a_m \cos(2m\pi\zeta) + b_m \sin(2m\pi\zeta)\right] e^{-4m^2\pi^2 \mathcal{D}_c t/L^2},
\qquad \tau_1 = \frac{L^2}{4\pi^2 \mathcal{D}_c},
\label{eq:hold}
\end{equation}
with $a_m$, $b_m$ set by the state at the end of the ramp. Integrating (\ref{eq:hold}) gives the displacement modes $\sin(2m\pi\zeta)$ and $\cos(2m\pi\zeta) - 1$, both zero at both plates.

\subsection{Compression, unloading}

At $t = 0$ the plates are withdrawn from a gel of uniform strain $\varepsilon_0 = \Delta L/L$. Both faces are stress free,
\begin{equation}
\varepsilon(0,t) = \varepsilon(L,t) = 0, \qquad \varepsilon(\zeta,0) = \varepsilon_0,
\label{eq:unloadbc}
\end{equation}
so
\begin{equation}
\varepsilon(\zeta,t) = \varepsilon_0 \sum_{k\ \mathrm{odd}} \frac{4}{k\pi} \sin(k\pi\zeta)\, e^{-k^2\pi^2 \mathcal{D}_c t/L^2},
\qquad \tau_1 = \frac{L^2}{\pi^2 \mathcal{D}_c},
\label{eq:unload}
\end{equation}
and the thickness recovers as
\begin{equation}
\frac{\delta L(t)}{\delta L(\infty)} = 1 - \sum_{k\ \mathrm{odd}} \frac{8}{(k\pi)^2}\, e^{-k^2\pi^2 \mathcal{D}_c t/L^2}.
\label{eq:unloadtrace}
\end{equation}
The displacement modes are $\cos(k\pi\zeta)$ in the frame of the compressed gel. The uniform mode of (\ref{eq:hold}) is absent here and the $m = 1$ rate of (\ref{eq:hold}) is four times the $k = 1$ rate of (\ref{eq:unload}): at equal $\mathcal{D}_c$ the unload is four times slower than the hold.

\subsection{Permeation, approach to steady state}

The network is pinned at the support, $u_z(0,t) = 0$, and the feed face is free. The pore pressure drop $\Delta P(t)$ is carried by the network at the support and vanishes at the feed face:
\begin{equation}
\varepsilon(0,t) = \frac{\Delta P(t)}{M}, \qquad \varepsilon(L,t) = 0, \qquad \varepsilon(\zeta,0) = 0.
\label{eq:permbc}
\end{equation}
The steady state of (\ref{eq:upde}) under (\ref{eq:permbc}) is the planar profile,
\begin{equation}
u_z(\zeta,\infty) = -\frac{\Delta P L}{M}\left(\zeta - \frac{\zeta^2}{2}\right), \qquad u_F \equiv u_z(L,\infty) = -\frac{\Delta P L}{2M},
\label{eq:permss}
\end{equation}
which is (S127) with $z \to -z$. Writing the applied history as increments $w_j$ at $t_j$, $\Delta P(t) = \Delta P \sum_{t_j < t} w_j$, $\sum_j w_j = 1$, the transient is
\begin{equation}
u_z(\zeta,t) = \sum_{k=1}^{\infty} A_k \left[1 - \cos(k\pi\zeta)\right] H_k(t), \qquad
H_k(t) = \sum_{t_j < t} w_j \left[1 - e^{-k^2\pi^2 \mathcal{D}_c (t - t_j)/L^2}\right],
\qquad \tau_1 = \frac{L^2}{\pi^2 \mathcal{D}_c},
\label{eq:perm}
\end{equation}
and the feed-face displacement follows
\begin{equation}
\frac{u_F(t)}{u_F(\infty)} = \sum_{k\ \mathrm{odd}} \frac{8}{(k\pi)^2}\, H_k(t).
\label{eq:permtrace}
\end{equation}

\subsection{Fitting}

In each case the mode amplitudes are free and $\mathcal{D}_c$ is the single nonlinear parameter, found by least squares over the displacement snapshots of the first $6\tau_1$ ($\tau_1$ of the fit itself, iterated), each mode decay averaged over the output window of the snapshot it is compared with. The thickness traces (\ref{eq:unloadtrace}) and (\ref{eq:permtrace}) give a second $\mathcal{D}_c$ with one amplitude, $\delta L(\infty)$ or $u_F(\infty)$. Then $M = \Delta P L/(2|u_F|)$ from (\ref{eq:permss}) and $\kappa = \mathcal{D}_c/M$.

\end{document}
